% !TeX root = ../教材题库_全部条目_转写工作稿.tex
% Source: XJ-B1-C06-EX-P218-Y1196
% 第6章 / 6.2 抽样 / 教材 p.218、219 / PDF 第225、226页
% OCR未校对：公式与文字均可能有误；未完成数学TeX转写；禁止直接发给学生
% 原题及相邻原解截图（核对用）：../题目原图/XJ-B1-C06-EX-P218-Y1196-01.webp, ../题目原图/XJ-B1-C06-EX-P218-Y1196-02.webp
% 本文件只在转写完成后作为题干使用；OCR 条目尚不是可靠题干。
\begin{bankitem}{6.2 例题 例题1}{第6章 / 6.2 抽样 / 教材 p.218、219 / PDF 第225、226页}{XJ-B1-C06-EX-P218-Y1196}{OCR未校对：公式与文字均可能有误；未完成数学TeX转写；禁止直接发给学生}
\begin{ocrdraft}
[教材p.218 / PDF第225页]
\par
00一家网络音乐平台就网络用户对某一特色栏目的喜爱程度进行调查，参
\par
与网络投票的总人数为50000．网络用户对栏目的评价如下表所示：
\par
星级评价
\par
票数
\par
9986
\par
15607
\par
14885
\par
6034
\par
3488
\par
该音乐平台为进一步了解用户的具体想法和意见，打算从上述50000人中抽取
\par
1佣人进行调查，应怎样进行抽样？
\par
分析因为观众对栏目喜爱程度差异较大，故应采用分层抽样．同时，每层随
\par
机抽取的数量比例，可与总体中各层票数所占比例相同．
\par
解采用分层抽样，其总体容量为50000．
\par
[教材p.219 / PDF第226页]
\par
9986
\par
9986
\par
“五星评价”占
\par
应抽取100\ensuremath{\times}
\par
20（人）；
\par
50000，
\par
50000
\par
15607
\par
15607
\par
“四星评价”占
\par
应抽取100\ensuremath{\times}
\par
31（人）；
\par
50000，
\par
50000
\par
三星评价”占
\par
14885
\par
应抽取100\ensuremath{\times}
\par
14885
\par
30（人）；
\par
50000，
\par
50000
\par
二星评价”占
\par
6034
\par
应抽取100\ensuremath{\times}
\par
6034
\par
12（人）；
\par
50000，
\par
50000
\par
3488
\par
3488
\par
一星评价”占
\par
应抽取100\ensuremath{\times}
\par
7（人）．
\par
50000，
\par
50000
\par
即按照“五星评价”至“一星评价”分别抽取20人，31人，30人，12人，
\par
从上面的抽样过程可以看出，分层抽样保证了样本中包含各种特征的抽样单位，
\par
样本的结构与总体的结构保持一致性，这对提高样本的代表性是非常重要的．在后
\par
面学习用样本估计总体时，我们将会了解到分层抽样既可以对总体特征进行估计，也
\par
可以对各层的特征进行估计．这些优点使分层抽样在实践中得到了广泛的应用．
\par
前面介绍了简单随机抽样和分层抽样．对一个实际问题而言，究竟用何种抽样
\par
方法应视具体情况而定．
\par
因为实际问题的复杂性，我们进行抽样调查时，除了要考虑样本的代表性外，
\par
还要考虑调查实施中的人力、物力等诸多因素，因此，许多时候并不是单一使用某
\par
种抽样方法，而是将多种抽样方法组合起来使用．
\end{ocrdraft}
\end{bankitem}
